\section{Quantifying Error of Unitary Operations}\label{sec:error}

\begin{remark}[Soundness of Definition of Error]
	Let $\ket{\psi}$ be the initial state of the system, which is evolved according to unitary $U$ or $V$, and then measured via POVM measurement $M$. 
	Let $P_U$ and $P_V$ denote the probability of obtaining the corresponding measurement outcome when the operator $U$ or $V$ are performed, respectively. 

	Then 
	$$
	|P_U - P_V| = | \bra{\psi} U^\dagger M U \ket{\psi} - \bra{\psi} V^\dagger M V \ket{\psi} |
	$$
	Let $\ket{\Delta} := (U - V)\ket{\psi}$, we have 
	\begin{align*}
		|P_U - P_V| 
		& = | \bra{\psi} U^\dagger M \ket{\Delta} + \bra{\Delta} M V \ket{\psi} | & \\
		&\leq | \bra{\psi} U^\dagger M \ket{\Delta} | + |\bra{\Delta} M V \ket{\psi} | &\\
		& \leq \| \ket{M^\dagger U \psi} \| \cdot \| \ket{\Delta} \| + \| \ket{M V \psi} \| \cdot \| \ket{\Delta} \|  & \text{(by Cauchy-Schwarz)}\\ 
		& \leq 2 \| \ket{\Delta} \| &  \| \ket{M^\dagger U \psi}\| \leq 1\\
		& = 2 E(U, V) &
	\end{align*}

	This means if $E(U, V)$ is small, the probability of obtaining a different measurement outcome is also small, which is exactly what we expect for the definition of error.
\end{remark}

\begin{proposition}[Cumulation of Error]\label{prop:cumulation-of-error}
	If we have a sequence of unitary operations $U_1, U_2, \cdots, U_n$ and their approximations $V_1, V_2, \cdots, V_n$ with error $E(U_i, V_i)$, then 
	$$
	E(U_1 U_2 \cdots U_n, V_1 V_2 \cdots V_n) \leq \sum_{i=1}^n E(U_i, V_i).
	$$
\end{proposition}

\begin{proof}
	Let us first consider $E(U_1U_2, V_1V_2)$
	By definition
\begin{align*}
	E(U_1U_2, V_1V_2) 
	& = \max_{\|\ket{\psi}\| = 1} \| (U_1 U_2 - V_1 V_2) \ket{\psi} \|&\\
	& = \max_{\|\ket{\psi}\| = 1} \| (U_1 U_2 - U_1 V_2 + U_1 V_2 - V_1 V_2) \ket{\psi} \| &\\
	& \leq \max_{\|\ket{\psi}\| = 1} \| \| (U_1 U_2 - U_1 V_2) \ket{\psi} \| + \| (U_1 V_2 - V_1 V_2) \ket{\psi} \| & \\
	& = E(U_2, V_2) + E(U_1, V_1) &
\end{align*}
The general case will follow from induction.
\end{proof}

The definition of error is echoes the definition of matrix 2-norm.

\begin{definition}[Matrix 2-norm]
	The matrix 2-norm of a matrix $M$ is defined as 
	$$
	\| M \| = \max_{\ket{\psi}} \frac{\| M \ket{\psi} \|}{\| \ket{\psi} \|}.
	$$
\end{definition}

\begin{proposition}[Computation of 2-Norm]
	Assuming $M$ is normal\footnote{A matrix $M$ is normal if $M M^\dagger = M^\dagger M$.}, 
	Then
	$$
	\| M \| = \max_{\ket{\psi}} \frac{\| M \ket{\psi} \|}{\| \ket{\psi} \|} = ||\lambda_{\max}|,
	$$
	Where $\lambda_{\max}$ is the eigenvalue of $U - V$ with the largest modulus.
\end{proposition}

\begin{proof}
	This is the standard result about matrix norm. 

	Since $M$ is normal by spectral theorem $M$ can be decompose into $\lambda_1\ket{u_1}\bra{u_1} + \lambda_2 \ket{u_2} \bra{u_2}$, where $\lambda_1, \lambda_2$ are its eigenvalues and $\ket{u_1}, \ket{u_2}$ are its orthogonal eigenvectors.

	Therefore any vector $\alpha$ can be written as $\alpha = a \ket{u_1} + b \ket{u_2}$ for some $a, b \in \mathbb{C}$, and we have

	\begin{align*}
		\frac{\|M \ket{\alpha}^2 \|}{\| \ket{\alpha}\|^2} 
				 &= \frac{\|M (a \ket{u_1} + b \ket{u_2})\|^2}{|a|^2 + |b|^2} 
				 &\\ 
				 & = \frac{\|a \lambda_1 \ket{u_1} + b \lambda_2 \ket{u_2}\|^2}{|a|^2 + |b|^2} 
				 &\\ 
				 & \leq \frac{ |a|^2 |\lambda_1|^2 + |b|^2 |\lambda_2|^2}{|a|^2 + |b|^2}
				 & \text{triangular inequality}\\ 
				 & \leq \frac{\max(|\lambda_1|, |\lambda_2|)^2 (|a|^2 + |b|^2)}{|a|^2 + |b|^2}
				 &\\ 
				 & = \max(|\lambda_1|, |\lambda_2|)^2 &
	\end{align*}
	Taking the square root on both sides concludes the proof.
\end{proof}

\begin{corollary}[Computation of Error]\label{cor:computation-of-error}
	Let $U$ and $V$ be two unitary operations, then 
	$$
	E(U, V) = \| U - V \| = \| UV^\dagger - I \| = |\lambda_{\max}|,
	$$
	where $\lambda_{\max}$ is the eigenvalue of $UV^\dagger - I$ with the largest modulus.
\end{corollary}

\begin{proof}
	$ \| U - V \| = \| UV^\dagger - I \| $ follows from proposition \ref{prop:cumulation-of-error} on cumulation of error. 
	$UV^\dagger - I$ is normal because 
	\begin{align*}
		(UV^\dagger - I)(UV^\dagger - I)^\dagger 
		& = (UV^\dagger - I)(VU^\dagger - I) &\\
		& = UV^\dagger VU^\dagger - UV^\dagger - VU^\dagger + I &\\
		& = I - UV^\dagger - VU^\dagger + I &\\
		& = (VU^\dagger - I)(UV^\dagger - I) = (UV^\dagger - I)^\dagger (UV^\dagger - I) &
	\end{align*}
\end{proof}
